\begin{frame}{Példa lefutás: második nem-fa él $(b, g)$ (1.)}
\begin{columns}[T]
  \begin{column}{0.46\textwidth}
    \small
    \textbf{Második nem-fa él $(b, g)$:}
    \begin{itemize}
      \item Meghívjuk PathLabel(b, g, (b, g))-t %% TODO check if this is right??
      % \item We get the RIGHT branch plan with k1 = OUT[g], k2 = IN[b] and so again k1 <k2
      \item RIGHT plan-t kapjuk meg. $k_1$ = OUT[g] (= 7), $k_2$ = IN[b] (= 10), szóval $k_1$ < $k_2$.
      % \item The traversal begins with b(line 11) and since b has not yet been visited then (b,e) gets the non-MST edge (b,g), where e is the parent of b.
      \item A bejárást b-vel kezdjük és mivel b-t még nem látogattuk meg ezért (b, e) él a (b, g) nem-MST élt kapja meg csereélként. ($R_{(b, e)} = (b, g)$)
      % \item The disjoint set is linked (line 15) so b’s disjoint set gets e’s label and the subpath (b,e) is compressed.
      \item b és e diszjunkt halmazait link-eljük. Ezzel tömörítjük a (b, e) részutat.
    \end{itemize}
  \end{column}
  \begin{column}{0.54\textwidth}
    \centering
    \begin{tikzpicture}[
      scale=1,
      vtx/.style={circle, draw=black, fill=white, inner sep=1pt, minimum size=16pt},
      mstedge/.style={line width=2.5pt, draw=black},
      otheredge/.style={line width=0.8pt, draw=black},
      font=\small
    ]

    % --- Node coordinates ---
    \node[vtx] (a) at (-4.2, 1) {a};
    \node[vtx] (b) at (-2.2, 2) {b};
    \node[vtx] (d) at (-3, -1.2) {d};
    \node[vtx] (e) at (0, 1) {e};
    \node[vtx] (f) at (-0.5, -1.2) {f};
    \node[vtx] (g) at (2, 2) {g};
    \node[vtx] (h) at (2.5, 0) {h};

    % -- root node
    \node[rootvertex] (c) at (-2.0, 0.2) {c};

    % --- MST edges (thick) ---
    \draw[mstedge] (d) -- (f)
      node[midway, below=2pt, fill=white, inner sep=1pt, font=\scriptsize]{1};
    \draw[mstedge] (g) -- (h)
      node[midway, right=2pt, fill=white, inner sep=1pt, font=\scriptsize]{3};
    \draw[mstedge] (e) -- (f)
      node[midway, right=2pt, fill=white, inner sep=1pt, font=\scriptsize]{4};
    \draw[mstedge] (b) -- (e)
      node[midway, below=2pt, fill=white, inner sep=1pt, font=\scriptsize]{5};
    \draw[mstedge] (c) -- (d)
      node[midway, left=3pt, fill=white, inner sep=1pt, font=\scriptsize]{7};
    \draw[mstedge] (a) -- (d)
      node[midway, left=3pt, fill=white, inner sep=1pt, font=\scriptsize]{9};
    \draw[mstedge] (e) -- (h)
      node[midway, above=2pt, fill=white, inner sep=1pt, font=\scriptsize]{2};

    % --- Non-tree edges (thin) ---
    \draw[otheredge, draw=red] (b) -- (g)
      node[midway, above=2pt, fill=white, inner sep=1pt, font=\scriptsize]{8};
    \draw[otheredge] (a) -- (b)
      node[midway, above=3pt, fill=white, inner sep=1pt, font=\scriptsize]{10};
    \draw[otheredge] (c) -- (e)
      node[midway, above=2pt, fill=white, inner sep=1pt, font=\scriptsize]{11};
    \draw[otheredge] (b) -- (c)
      node[midway, left=2pt, fill=white, inner sep=1pt, font=\scriptsize]{12};
    \draw[otheredge] (f) -- (h)
      node[midway, above=2pt, fill=white, inner sep=1pt, font=\scriptsize]{13};
    % aktuálisan vizsgált nem-fa él
    \draw[otheredge, line width=1.2pt] (e) -- (g)
      node[midway, above=2pt, fill=white, inner sep=1pt, font=\scriptsize]{6};

    % --- DFS (IN, OUT) pairs ---
    \node[font=\scriptsize] at ($(a)+(-0.2,0.5)$) {(13, 14)};
    \node[font=\scriptsize] at ($(b)+(0,0.6)$) {(9, 10)};
    \node[font=\scriptsize] at ($(c)+(-0.6,0.4)$) {(1, 16)};
    \node[font=\scriptsize] at ($(d)+(-0.1,-0.6)$) {(2, 15)};
    \node[font=\scriptsize] at ($(e)+(0,0.5)$) {(4, 11)};
    \node[font=\scriptsize] at ($(f)+(0.0,-0.6)$) {(3, 12)};
    \node[font=\scriptsize] at ($(g)+(0,0.6)$) {(6, 7)};
    \node[font=\scriptsize] at ($(h)+(0.4,0.5)$) {(5, 8)};

    \end{tikzpicture}
  \end{column}
\end{columns}
\end{frame}